\scanpage{7}
\uline{Generalization over an arbitrary base.}

A \uline{standard curve} over $S$ (schematic multiplicity, say) will be defined constructively in terms of a system a) b) c) \inserted{as} above, i.e.
\[
(5)\quad\left\{\begin{aligned}
&\text{a) }Y,\text{ proper and smooth scheme }\inserted{\text{over }S}\text{ everywhere of}\\
&\text{relative dimension }1;\\
&\text{b) }\widetilde A\text{ and }I,\text{ two subschemes, of }Y,\\
&\text{finite étale over }S,\quad\deleted{\unsure{\text{with}|\text{and}}}\ \widetilde A\cap I=\unsure{\phi|\varnothing};\\
&\text{c) }\sigma_{\widetilde A}\text{ an involution of }\widetilde A,\text{ acting}\\
&\text{freely.}
\end{aligned}\right.
\]
One then constructs $\widehat X=\deleted{X}\inserted{Y}/\sigma$, containing $A\simeq\widetilde A/\sigma$ and $I$ as \inserted{closed} \unsure{subschemes|subschemes} finite étale over $S$, and one defines $X=\widehat X\setminus I$. One can show that the functor
\[
(6)\qquad(Y,\widetilde A,\inserted{I,}\sigma_{\widetilde A})\longmapsto X
\]
\margin{\uline{NB} NB this functor is compatible with inverse images\ldots}
from systems (5) (for isomorphisms) to relative \unsure{schemes|schemes} \inserted{(for isomorphisms)}, is \uline{fully faithful}.
\margin{false as it stands!}
\mspage{6}
